\chapter{Anti-Relaxation Coatings for Alkali-Vapour Cell}
\label{app:anti_relaxation_coatings}
Anti-relaxation coatings are used to reduce the loss of atomic polarization and ground-state coherence caused by collisions with the inner surface of an alkali-vapour cell. Their function is not to prevent atoms from reaching the wall, but to modify the atom--surface interaction so that a wall encounter produces a smaller perturbation of the internal state. The coating therefore acts as an effective boundary material between the alkali vapour and the underlying glass substrate. Its performance depends on both the chemical structure of the coating and the microscopic properties of the deposited film.

Following the classification framework established by Chi et al.~\cite{Aiacoavc_ChiEtAl_2020}, the principal anti-relaxation coating families for alkali-vapour cells include straight-chain alkanes, alkenes, and organochlorosilanes. These materials differ in their chemical bonding, thermal stability, and anti-relaxation performance, leading to different compromises between spin preservation and operating temperature.

Straight-chain alkanes, such as paraffin, form ``wet'' films mainly through physical adsorption on the inner glass surface. These coatings were among the first successful anti-relaxation coatings and can allow alkali atoms to undergo many wall collisions before depolarization, with typical reported values around $N_{\mathrm{b}}\sim10^{4}$. Their chemical structure is dominated by saturated C--C bonds. These bonds have relatively low polarizability and weak adsorption energies, which helps reduce spin-randomizing interactions during wall contact. However, paraffin-based coatings are limited by thermal stability and are generally unsuitable for high-temperature operation.

Alkene coatings are also physically adsorbed films, but they contain unsaturated C=C bonds. These bonds are more polarizable than saturated C--C bonds, and early intuition might suggest that this would be harmful for spin preservation. However, experiments discussed in the coating literature show that C=C bonds are not necessarily strongly depolarizing, and alkene coatings can reach very large bounce numbers, in some cases approaching $N_{\mathrm{b}}\sim10^{6}$~\cite{Aiacoavc_ChiEtAl_2020}. Their principal limitation is again temperature stability, since their anti-relaxation performance can degrade at temperatures relevant for high-density alkali vapours. This makes them attractive from the point of view of spin survival, but less suitable when high optical depth requires elevated cell temperature.

Organochlorosilanes, such as octadecyltrichlorosilane (OTS), represent a different chemical strategy. Instead of forming only a physically adsorbed wet film, these coatings can form chemically bonded dry films. The chlorosilane end groups react with hydroxyl groups on the glass or silicon surface, producing siloxane bonds that anchor the organic layer to the substrate~\cite{Aiacoavc_ChiEtAl_2020}. This chemical attachment improves thermal robustness and allows operation at substantially higher temperatures than typical paraffin or alkene coatings. The trade-off is that the spin survival is usually lower, with reported maximum values closer to $N_{\mathrm{b}}\sim10^{3}$ than to the best alkene coatings. Thus, organochlorosilanes are not necessarily the best coatings in terms of bounce number, but they remain important for compact warm-vapour memories because they can operate in temperature regimes where high optical depth becomes accessible.

The coating literature also shows that the microscopic anti-relaxation mechanism is not yet completely understood. For paraffin, models based on adsorption, electron-randomization relaxation, spin-exchange relaxation, and relaxation associated with atoms in stems or reservoirs have been proposed. During a wall encounter, an atom may remain adsorbed on the surface for a finite dwell time before desorbing back into the vapour. While adsorbed, it can experience local magnetic fields, electric-field gradients, spin-orbit interactions, or interactions with nuclear spins in the substrate, all of which may contribute to random phase accumulation and spin relaxation~\cite{OP_Happer_1972,Aiacoavc_ChiEtAl_2020}. The relative importance of these processes depends on the coating material and operating conditions.

Chi et al.~\cite{Aiacoavc_ChiEtAl_2020} further emphasize that surface-science techniques are increasingly used to study coating composition and morphology. Common characterization methods include atomic force microscopy, ellipsometry, X-ray diffraction, and contact-angle measurements. These techniques provide information about film thickness, roughness, molecular ordering, surface coverage, and other structural properties that influence the interaction between the alkali atom and the coated surface.

Consequently, coating performance depends not only on the chemical family, but also on the fabrication procedure, thickness, roughness, surface coverage, molecular ordering, adsorption energy, alkali exposure, and operating temperature. Quantities such as $N_{\mathrm{b}}$, atomic dwell time, hyperfine relaxation, and memory storage time should therefore be regarded as effective experimental characterizations of a complex microscopic surface rather than as intrinsic material constants.
